\begin{abstract}
Post-hoc model merging has emerged as an efficient, training-free paradigm for multi-task learning, enabling separate task-specific expert models independently fine-tuned from a shared initialization to be directly composed in weight space. To mitigate task conflicts, recent frameworks employ test-time adaptation (TTA) to dynamically optimize layer-wise merging coefficients on small calibration streams. However, we identify a critical, overlooked vulnerability in these adaptive paradigms, termed \textbf{Quantization-Operator Overfitting}: unconstrained test-time coefficient optimization converges to extremely sharp local minima that collapse under downstream deployment post-training quantization (PTQ). 

To resolve this bottleneck, we present \textbf{CR-PolySACM (Clipping-Regularized Sharpness-Aware Subspace Model Merging)}, a unified framework that combines global structural constraints with local landscape flatness optimization. First, we identify a fundamental \emph{task-vector norm scale pathology} in standard sharpness-aware optimization, where a 50-fold discrepancy in layer-wise task-vector norms renders the optimizer blind to highly sensitive, low-norm layers such as final layer normalization, triggering optimization instability and collapse. Second, we introduce \textbf{Clipping-Regularized SACM (CR-SACM)}, which clips task-vector norms to a robust minimum floor to resolve scale-blindness without triggering gradient explosion. By integrating CR-SACM with a low-degree depth-dependent polynomial subspace parameterization, our proposed CR-PolySACM successfully maps the optimization trajectory into well-conditioned, robust regions of the landscape. We refer to our proposed framework interchangeably as \textbf{CR-PolySACM} and the shorthand \textbf{PolySACM} throughout this work. 

Our empirical evaluation across six diverse hardware-relevant quantization schemas (ranging from FP32 down to INT4 symmetric per-channel) demonstrates that CR-PolySACM consistently stabilizes model composition, achieving state-of-the-art joint multi-task accuracies. Under aggressive 4-bit quantization, CR-PolySACM achieves a joint mean accuracy of \textbf{19.07\%}, outperforming the previous state-of-the-art PolyMerge baseline (\textbf{18.10\%}) by nearly \textbf{+1.0\%} and demonstrating the critical synergy of global structural subspaces and local flatness regularization for robust edge deployment.
\end{abstract}
